
In this paper, I provide a novel model framework to analyze the effects of immigration on business activity in the United States. I contribute to a rich literature of empirical work by performing my analysis in a more granular setting and with frontier identification methods in the study of migration choice. I provide theoretical and empirical evidence that migrants in the US are positively selected on managerial talent and they can cause increases in firm size. However, this study is situated in a capital-constrained setting, as shown, much of the underlying talent ability is depressed by a lack of capital. An interesting extension of this paper would be examining whether the capital constraints affect immigrants and natives differently or if the groups have differential access to capital sources. These findings should be interpreted in the context of the study, which is U.S. counties from 1975-2010. The United States has a unique immigration system and business environment that attracts particular types of individuals \footnote{See \cite{ImmigrationinAmericanEconomicHistory} for a comprehensive history of immigration to the United States}. While my IV approach deals with the endogeneity of where an immigrant settles within the US, it does not account for the selection of who chooses to become an immigrant. Making the decision to move to a foreign country carries a great amount of risk, and integrating into a new society poses its own unique challenges. As such, it is not surprising that the choice of migration and managerial talent is correlated. In other words, selecting into becoming a migrant and having high managerial talent go hand in hand. Another strand of interesting future work would be to study who chooses to become an immigrant.

However, migration is not always a choice. Sometimes domestic factors force individuals to move to a foreign country, this can include wars, famines, or extreme weather events. This is relevant as all types of immigration are not the same. Studies on immigration must understand and make the distinction between different types of migrants. The positive selection of migrants is not a necessity and can differ from place to place and over time periods. 

In conclusion, I hope this paper contributes to the growing literature on the unique study of migration, along with its many unique effects. I also provide a version of the Span-of-Control model which can easily be adapted to study any two groups with different managerial talents, such as male vs. female, educated vs. non-educated, or young vs. old manager and not just immigrants versus native entrepreneurs.